\begin{abstract}

Human population structure has been characterized primarily through single nucleotide variants (SNVs), yet the genome harbors millions of insertions/deletions (INDELs) and thousands of structural variants (SVs) whose population-genetic behavior has not been systematically compared. Here, using the 2022 high-coverage 1000 Genomes Project panel (3,202 individuals, 26 populations), we performed the first unified comparison of population structure, differentiation, and allele frequency dynamics across SNVs (12.7M), INDELs (3.3M), and SVs (18.5K). All three variant types recovered the same continental population clusters, and pairwise $F_\mathrm{ST}$ values were highly correlated across types ($r > 0.997$). However, mean $F_\mathrm{ST}$ followed a systematic gradient---SNV (0.080) $>$ INDEL (0.072) $>$ SV (0.059)---indicating that population differentiation is attenuated for larger variants. A convergence analysis, in which SNVs were subsampled at counts from 500 to 1,000,000, revealed that SNV-INDEL Procrustes concordance converged to $r \approx 0.99$ while SNV-SV concordance plateaued at $r \approx 0.60$ regardless of variant count, establishing the SV discordance as irreducible biological signal rather than a statistical artifact. Size-stratified analysis across two orders of magnitude (1--200~bp) demonstrated a continuous, monotonic decline in $F_\mathrm{ST}$ with variant size, consistent with stronger purifying selection on larger genomic rearrangements. Insertions showed 25\% lower $F_\mathrm{ST}$ than deletions of comparable size ($F_\mathrm{ST}$ = 0.046 vs.\ 0.062), a pattern consistent with asymmetric selective constraint, though differential short-read ascertainment cannot be excluded. These findings reveal a size-dependent selection gradient as a previously unrecognized dimension of human population genomic variation and provide a convergence analysis framework for distinguishing power effects from genuine biological discordance in cross-type comparisons.

\end{abstract}

\noindent\textbf{Author summary}\quad
Population genetics studies almost exclusively use single-letter DNA changes (SNVs) to characterize the genetic relationships among human populations. But our genomes also carry millions of small insertions and deletions and thousands of larger structural rearrangements. We asked whether these different types of genetic variation tell the same story about human population structure. Analyzing over 3,000 individuals from the 1000 Genomes Project, we found that all variant types identify the same major population groups, but larger variants show systematically weaker differences between populations---structural variants are about 28\% less differentiated than SNVs. By varying the number of variants used in our comparisons, we demonstrated that this discordance is not simply due to having fewer structural variants to analyze; it reflects a genuine biological difference. Differentiation decreases smoothly and continuously as variant size increases from 1~bp to over 200~bp, consistent with stronger natural selection against larger genomic changes. We also discovered that DNA insertions are under noticeably stronger selective pressure than deletions of similar size. These results establish that variant type matters for population genetics and provide practical guidance for studies that use non-SNV variants.
